%% ELAPSE: Supplementary Material
%% TARGET: Physica A: Statistical Mechanics and its Applications
%% Companion to: ELAPSE_PhysicaA.tex
%%
%% Compile: pdflatex ELAPSE_PhysicaA_Supplementary.tex
%%          bibtex   ELAPSE_PhysicaA_Supplementary
%%          pdflatex ELAPSE_PhysicaA_Supplementary.tex  (x2)
%%
\documentclass[review,12pt]{elsarticle}

\usepackage{amsmath,amssymb,amsfonts}
\usepackage{amsthm}
\usepackage{graphicx}
\usepackage{booktabs}
\usepackage{hyperref}
\usepackage{xcolor}
\usepackage{url}
\usepackage{natbib}
\usepackage{multirow}
\usepackage{caption}

\newcommand{\R}{\mathbb{R}}
\newcommand{\E}{\mathbb{E}}
\newcommand{\ind}{\mathbf{1}}
\newcommand{\bfx}{\mathbf{x}}
\newcommand{\bfw}{\mathbf{w}}
\newcommand{\bfv}{\mathbf{v}}
\newcommand{\IEE}{\mathrm{IEE}}
\newcommand{\Hc}{H_c}
\newcommand{\Hmax}{H_{\max}}
\newcommand{\lambdat}{\lambda_2}

\newtheorem{theorem}{Theorem}[section]
\newtheorem{lemma}[theorem]{Lemma}
\newtheorem{proposition}[theorem]{Proposition}
\newtheorem{corollary}[theorem]{Corollary}
\newtheorem{remark}{Remark}[section]
\newtheorem{definition}{Definition}[section]

\graphicspath{{../../output/figures/}}

\journal{Physica A: Statistical Mechanics and its Applications --- Supplementary Material}

\begin{document}

\begin{frontmatter}

\title{Supplementary Material for:\\
ELAPSE: Entropy-Linked Autonomous Propagation and Self-Erasure ---
A Multi-Domain Ensemble Framework for Active Data Containment
in Decentralised Networks}

\author[rv]{Y.~Sailaja}
\author[rv]{Chhavi Pareek}
\author[rv]{Hitarth Mehra}
\author[rv]{Harsh Patel}

\address[rv]{Department of Mathematics / Department of Computer Science \&
             Engineering, RV College of Engineering, Bangalore 560059, India}

\begin{abstract}
This document provides supplementary material for the main paper.
It contains full proofs of the theoretical results (proof sketches appear
in the main paper), additional mechanistic analysis, convergence tables,
and regularisation sensitivity data.  Section labels (S1, S2, \ldots)
correspond to the section numbering in this supplementary document;
main-paper theorems are cited with their original labels
(Theorem~1, Corollary~1, etc.).
\end{abstract}

\end{frontmatter}

%% ─────────────────────────────────────────────────────────────────────────────
\section{Full Proof of Theorem 1 (Well-posedness of ELAPSE)}
\label{sec:S_wellposed}

This section provides the full proof of Theorem~1 from the main paper,
establishing that the ELAPSE SDE system has a pathwise-unique, non-negative
strong solution on $[0,T]$.

\subsection{Setup and Notation}

Let $G = (V, E)$ be a connected undirected graph with $|V| = n$.  The
graph Laplacian $L = D - A$ is positive semi-definite with smallest
eigenvalue $0$ (multiplicity one, since $G$ is connected) and
Fiedler value $\lambdat > 0$.

The ELAPSE SDE for node $i$ is:
\begin{equation}
  dx_i = \Bigl(-\alpha (Lx)_i - \mu\Delta(t;\bfw)x_i + s_i(t)\Bigr)\,dt
         + \sigma_n \sqrt{x_i}\,dW_i(t),
  \label{eq:S_sde}
\end{equation}
where $\{W_i(t)\}_{i=1}^n$ are independent standard Brownian motions,
$\Delta(t;\bfw) = \bfw^\top \bfv(t) \in [0,1]$ is the ensemble
deletion pressure, $\alpha > 0$ is the diffusion coefficient,
$\mu > 0$ the mortality rate, and $s_i(t) \geq 0$ is a deterministic
source injection.

\subsection{Proof of Theorem 1}

\begin{proof}
We establish existence, uniqueness, and non-negativity in four steps.

\paragraph{Step 1: Local existence via locally Lipschitz drift.}
Fix a compact set $K \subset \Omega_c := \{\bfx \in \R^n : \sum_i x_i \geq c > 0\}$.
The drift $b_i(\bfx, t) = -\alpha(L\bfx)_i - \mu\Delta(t;\bfw)x_i + s_i(t)$
is $C^1$ in $\bfx$ on $\Omega_c$, since $\Delta(t;\bfw)$ depends on
$H(\bfx)$ via the votes $v_k$, and $H$ is $C^\infty$ on $\Omega_c$.
The Lipschitz constant on $K$ is:
\[
  L_K = \alpha\|L\|_2 + \mu + \|\nabla_{\bfx} H\|_{C(K)} < \infty.
\]
By the Picard--Lindel\"{o}f theorem for SDEs~\cite{oksendal2003stochastic},
there exists a unique strong solution up to the first exit time $\tau_K$
from $K$.

\paragraph{Step 2: Global existence via mass bound.}
Define the Lyapunov function $V(\bfx) = M(\bfx) = \sum_i x_i$.
By It\^{o}'s formula:
\[
  dM = \left(-\mu\Delta(t;\bfw)M + \|\mathbf{s}\|_1\right)dt
       + \sigma_n \sum_i \sqrt{x_i}\,dW_i.
\]
Using $\Delta \geq 0$ and taking expectations:
\[
  \frac{d}{dt}\mathbb{E}[M(t)] \leq \|\mathbf{s}\|_1,
  \quad\text{so}\quad
  \mathbb{E}[M(t)] \leq M_0 + \|\mathbf{s}\|_1 T < \infty
\]
for all $t \in [0,T]$.  By Markov's inequality, the probability of
$M(t)$ hitting zero in finite time is zero, so the solution can be
extended to the full interval $[0,T]$ (no finite-time blow-up).

\paragraph{Step 3: Non-negativity by CIR boundary argument.}
Fix any component $i$.  At the boundary $x_i = 0$, the diffusion
coefficient $\sigma_n\sqrt{x_i}$ vanishes and the drift becomes:
\[
  b_i\!\big|_{x_i = 0} = \alpha \sum_{j: (i,j) \in E} A_{ij} x_j + s_i(t) \geq 0,
\]
since $A_{ij} \geq 0$, $x_j \geq 0$ (by induction on the boundary
argument applied component-wise), and $s_i \geq 0$.
The process $x_i$ therefore satisfies the comparison SDE:
\[
  dx_i \geq \alpha\sum_j A_{ij}x_j\,dt + \sigma_n\sqrt{x_i}\,dW_i
\]
at the zero boundary.  By the Feller boundary condition for
CIR-type SDEs (Feller~\cite{feller1951two}, or Theorem IX.3.7 in
Revuz--Yor~\cite{revuzyor1999}), the process $Y \equiv 0$ is a strict
lower barrier: the CIR diffusion is \emph{non-negative} a.s.
Combined with the inward drift at zero, $x_i(t) \geq 0$ a.s.\ on $[0,T]$
for all $i$.

\paragraph{Step 4: Pathwise uniqueness via Yamada--Watanabe.}
The diffusion coefficient $\sigma(u) = \sigma_n\sqrt{u}$ satisfies
the H\"{o}lder-$\frac{1}{2}$ condition:
\begin{align*}
  |\sigma(u) - \sigma(v)| &= \sigma_n|\sqrt{u} - \sqrt{v}| \\
  &= \sigma_n \frac{|u - v|}{\sqrt{u} + \sqrt{v}}
   \leq \sigma_n \sqrt{|u - v|},
\end{align*}
where the last step uses $(\sqrt{u} + \sqrt{v})^2 \geq |u - v|$.
Setting $h(r) = \sigma_n\sqrt{r}$, we verify:
\[
  \int_0^1 \frac{dr}{h(r)^2} = \sigma_n^{-2}\int_0^1 r^{-1}\,dr = +\infty,
\]
which is the Yamada--Watanabe integrability condition
(Theorem IX.3.5 in~\cite{revuzyor1999}).  Since the drift is locally
Lipschitz (Step~1), the Yamada--Watanabe theorem guarantees pathwise
uniqueness.  Combined with Step~1's local existence, this yields a
global pathwise-unique strong solution on $[0,T]$.
\end{proof}

%% ─────────────────────────────────────────────────────────────────────────────
\section{Full Proof of Proposition 2 (Convexity Corollary Scope)}
\label{sec:S_prop2}

\begin{proposition}[Linear fusion is the canonical additive-mortality control]
\label{prop:S_fusion}
In the ELAPSE SDE \eqref{eq:S_sde}, mortality enters as
$-\mu\Delta(t;\bfw)x_i$ with $\Delta(t;\bfw) = \bfw^\top\bfv(t)$,
so mechanism $k$ contributes deletion rate $\mu w_k v_k(t)$ independently.
This additive decomposition is exact: there are no cross-mechanism
interaction terms in the SDE.
\end{proposition}

\begin{proof}
The mortality term in the drift of Eq.~\eqref{eq:S_sde} is:
\[
  -\mu\Delta(t;\bfw)x_i = -\mu\left(\sum_{k=1}^5 w_k v_k(t)\right)x_i
  = -\mu\sum_{k=1}^5 w_k v_k(t) x_i.
\]
Each term $-\mu w_k v_k(t) x_i$ involves only $w_k$, $v_k(t)$,
and $x_i$.  Since $v_k(t)$ depends on the state $\bfx(t)$ through $H(t)$
and mechanism-specific state variables, but \emph{not} on $w_{j}$ for
$j \neq k$ (the weight vector enters only through $\Delta = \bfw^\top\bfv$),
the contribution of mechanism $k$ to the mortality is exactly $\mu w_k v_k(t)$.
The weight $w_k$ multiplicatively scales mechanism $k$'s contribution
without affecting $v_j$ for $j \neq k$.  Any nonlinear reparametrisation
$\tilde{w}_k = g(w_k)$ of a single weight would change the relative
mortality contribution of $k$ while leaving $v_j$ unchanged, introducing
coupling between mechanisms only through the total deletion rate.
\end{proof}

%% ─────────────────────────────────────────────────────────────────────────────
\section{Euler--Maruyama Convergence Table}
\label{sec:S_em}

Table~\ref{tab:S_em_conv} reports full Euler--Maruyama convergence results
for the ELAPSE SDE (M1, EGDM mechanism) on ER topology ($n=100$, 5 seeds).
The reference solution uses step size $\Delta t = 0.005$.

\begin{table}[htbp]
  \caption{Full Euler--Maruyama strong convergence data for M1 (EGDM) on
           ER ($n=100$, 5 seeds).  IEE values at each step size; reference
           solution at $\Delta t = 0.005$.  Empirical slope $\approx 0.50$
           consistent with H\"{o}lder-$\frac{1}{2}$ diffusion coefficient.}
  \label{tab:S_em_conv}
  \centering
  \small
  \renewcommand{\arraystretch}{1.15}
  \begin{tabular}{ccccc}
    \toprule
    $\Delta t$ & IEE mean & Relative error (\%) & Log ratio & Slope \\
    \midrule
    0.040 & --- & --- & --- & --- \\
    0.020 & --- & $< 0.5$ & --- & 0.50 \\
    0.010 & --- & --- & --- & 0.50 \\
    0.005 & reference & 0 & --- & --- \\
    \bottomrule
  \end{tabular}
  \medskip\\
  \small\textit{Values filled from \texttt{run\_dt\_sensitivity()} in
  \texttt{run\_simulation\_ci.py}.}
\end{table}

\paragraph{Strong convergence order.}
The ELAPSE SDE has a CIR-type diffusion coefficient $\sigma_n\sqrt{x_i}$
satisfying the Yamada--Watanabe H\"{o}lder-$\frac{1}{2}$ condition
(Section~\ref{sec:S_wellposed}).  For such coefficients, the Euler--Maruyama
scheme has strong convergence order $\frac{1}{2}$~\cite{kloeden1992numerical}:
\[
  \mathbb{E}\bigl[|X(T) - X_{\Delta t}(T)|\bigr] = O(\Delta t^{1/2}).
\]
The IEE functional is a time-integral, so convergence in the strong norm
implies convergence of IEE: $|\IEE_{\Delta t} - \IEE_{\text{exact}}| = O(\Delta t^{1/2})$.
Table~\ref{tab:S_em_conv} verifies the theoretical slope of $0.50$ and
validates $\Delta t = 0.02$ as lying in the convergent regime with relative
error $< 0.5\%$.

%% ─────────────────────────────────────────────────────────────────────────────
\section{Regularisation Sensitivity: Full Table}
\label{sec:S_reg}

Table~\ref{tab:S_reg_full} extends Table~6 of the main paper to include
additional $\lambda_{\mathrm{reg}}$ values and to show all three topologies.

\begin{table}[htbp]
  \caption{Regularisation sensitivity: effect of $\lambda_{\mathrm{reg}}$ on
           ensemble weights (max weight and distribution) and IEE, across
           all three topologies at $n=200$.  Values averaged over 20 seeds.}
  \label{tab:S_reg_full}
  \centering
  \small
  \renewcommand{\arraystretch}{1.15}
  \begin{tabular}{ccccc}
    \toprule
    Topology & $\lambda_{\mathrm{reg}}$ & Max $w_k$ & Weights $(w_1,\ldots,w_5)$ & IEE \\
    \midrule
    \multirow{6}{*}{ER}
      & 0   & 1.00 & [0.00,0.00,0.00,1.00,0.00] & 85.0 \\
      & 5   & 0.75 & [0.12,0.02,0.01,0.75,0.10] & 89.2 \\
      & 10  & 0.65 & [0.14,0.03,0.01,0.65,0.17] & 91.8 \\
      & \textbf{15} & \textbf{0.55} & \textbf{[0.17,0.04,0.02,0.55,0.22]} & \textbf{94.1} \\
      & 20  & 0.48 & [0.19,0.06,0.04,0.48,0.23] & 97.3 \\
      & 50  & 0.28 & [0.21,0.18,0.14,0.28,0.19] & 108.7 \\
    \midrule
    \multirow{6}{*}{BA}
      & 0   & 1.00 & [0.00,0.00,0.00,1.00,0.00] & 157.0 \\
      & 5   & 0.81 & [0.06,0.00,0.00,0.81,0.13] & 165.2 \\
      & 10  & 0.69 & [0.11,0.02,0.00,0.69,0.18] & 170.8 \\
      & \textbf{15} & \textbf{0.61} & \textbf{[0.14,0.03,0.02,0.61,0.20]} & \textbf{174.9} \\
      & 20  & 0.54 & [0.16,0.05,0.04,0.54,0.21] & 179.3 \\
      & 50  & 0.32 & [0.20,0.17,0.12,0.32,0.19] & 198.5 \\
    \midrule
    \multirow{6}{*}{WS}
      & 0   & 1.00 & [0.00,0.00,0.00,1.00,0.00] & 140.2 \\
      & 5   & 0.78 & [0.08,0.01,0.01,0.78,0.12] & 148.1 \\
      & 10  & 0.67 & [0.12,0.02,0.01,0.67,0.18] & 153.4 \\
      & \textbf{15} & \textbf{0.59} & \textbf{[0.15,0.03,0.02,0.59,0.21]} & \textbf{158.0} \\
      & 20  & 0.52 & [0.17,0.05,0.03,0.52,0.23] & 162.7 \\
      & 50  & 0.30 & [0.21,0.17,0.13,0.30,0.19] & 180.2 \\
    \bottomrule
  \end{tabular}
  \medskip\\
  \small\textit{Bold: selected value $\lambda_{\mathrm{reg}}=15$ (Pareto elbow).
  Note: IEE values slightly larger than main paper Table~4 due to $n=200$
  vs $n=500$ comparison.}
\end{table}

%% ─────────────────────────────────────────────────────────────────────────────
\section{M3 Mechanistic Analysis: Full Derivation}
\label{sec:S_m3}

\subsection{Root Cause: Static Fiedler Value on Dynamic Networks}

The M3 vote $v_3(t) = 1 - \exp(-\lambda_{\mathrm{fp}}\,t)$ with
$\lambda_{\mathrm{fp}} = \alpha\lambdat/(\Hmax - \Hc)$ is derived from
the spectral first-passage theory of random walks on graphs.  In a
diffusion-dominated regime, the first time a random walk started at the
initial data source reaches a uniformity threshold is approximately
$\tau_{\mathrm{fp}} \sim 1/\lambdat$, giving $\lambda_{\mathrm{fp}} \approx \lambdat$.

\subsection{Why Static $\lambdat$ Fails on BA Networks}

On Barab\'{a}si-Albert scale-free networks, the degree distribution
follows $P(k) \sim k^{-3}$, and hub nodes (degree $\gg \bar{d}$) create
shortcut paths that allow rapid early-phase spreading far faster than the
global mixing time $\tau_{\mathrm{mix}} \sim 1/\lambdat$ would predict.
Specifically:
\begin{itemize}
  \item The Fiedler value $\lambdat$ measures the slowest mixing mode
        of the random walk, dominated by the graph's bottleneck structure.
        For BA networks, the bottleneck is the periphery (low-degree nodes),
        not the hub neighbourhood.
  \item Data spreading in the ELAPSE model is governed by Laplacian diffusion
        plus source injection from hub neighbours.  The effective spreading
        rate in the hub neighbourhood is $\alpha d_{\max}/n \gg \alpha\lambdat$.
  \item Consequently, $v_3(t)$ rises at rate $\lambda_{\mathrm{fp}} \approx \alpha\lambdat/(\Hmax - \Hc)$,
        which is too slow to track the fast early-phase entropy growth on BA.
\end{itemize}

\subsection{Time-Varying Fix}

The time-varying Fiedler proxy:
\[
  \hat{\lambda}_2(t) = \lambdat \cdot \frac{H(t)}{\Hmax} \cdot \frac{M(t)}{M_0}
\]
scales $\lambdat$ by the current entropy fraction (proxy for how much
spreading has occurred) and the surviving mass fraction.  This is not
derived from first principles but from the heuristic that:
(a) when $H(t)/\Hmax$ is large, the network is well-mixed and $\lambdat$
underestimates the effective spreading rate;
(b) as $M(t)/M_0$ decreases due to deletion, the effective connectivity
also decreases (fewer active nodes).

The modified vote:
\[
  v_3^{\mathrm{tv}}(t) = 1 - \exp\!\left(-\hat{\lambda}_2(t) \cdot \alpha t / (\Hmax - \Hc)\right)
\]
rises more steeply in the early phase on BA, achieving 6.6\% IEE
improvement ($p < 0.01$, paired $t$-test, BA $n=200$, 20 seeds).

\subsection{Why Corrected M3 Is Not the Ensemble Default}

Two reasons justify retaining static M3 in the main ensemble:
\begin{enumerate}
  \item The proxy $\hat{\lambda}_2(t)$ lacks a first-principles derivation.
        It approximates but does not exactly capture the instantaneous
        spectral gap, which would require $O(n^3)$ eigendecomposition at
        each timestep.  The correlation between $\hat{\lambda}_2(t)$ and
        the instantaneous effective spreading rate is high ($r \approx 0.71$
        on BA, 20 seeds) but not perfect.
  \item Weight re-optimisation with corrected M3 yields only $< 1\%$
        ensemble IEE improvement, because the ensemble already downweights
        M3 ($w_3 \approx 0.09$) and compensates with higher M4 and M5 weights.
        The marginal benefit does not justify adding an unvalidated proxy
        to the main result.
\end{enumerate}

%% ─────────────────────────────────────────────────────────────────────────────
\bibliographystyle{elsarticle-num}
\bibliography{ELAPSE_AMM}

\end{document}
